% Derivation strategies — each is a subsection. Treat this section as
% the technical core of the paper: the reader should be able to
% understand WHY each strategy exists and WHEN it wins, not just what
% it does. Where a strategy has been published before (Strassen,
% Pan, DIS09 §3), cite cleanly; where it's our combination of
% existing ideas (the recombination + analytical mask search), say
% so.

A \emph{derivation} strategy takes one or more catalog entries and
produces a new bilinear algorithm at a target shape --- by combining
them upward (\emph{composition}: Kronecker, concatenation), splitting a
larger target onto a smaller scheme (\emph{decomposition}, projection),
or locally reducing the product set (peel, pair fusion). The strategies
below cover what the catalog's search and materialiser currently know how
to do; each has a different cost structure and a different range of
shapes where it wins.

We adopt the convention that the \emph{outer} scheme is the one whose
bilinear structure is reused, and the \emph{sub-schemes} (or
\emph{leaves}) are the ones used to compute each of the outer
scheme's bilinear products at a smaller shape.

\paragraph{A unifying view, and its ceiling: smoothing \(\omega\) across shapes.}
Almost every strategy below is \emph{local} in the same sense: it synthesises a
scheme at a target shape out of schemes at \emph{surrounding} shapes --- the
sub-blocks of a decomposition, the factors of a Kronecker product, the
concatenands of an axis split. Read on the landscape of \(\omega\) over shapes,
these strategies \emph{interpolate}: the rank they produce at \nmpshape{n}{m}{p}
is an arithmetic combination (sum, product) of the ranks --- hence the
\(\omega\) --- of the shapes they consume, so the new local \(\omega\) is a
blend of its neighbours'. This has a structural ceiling. A pure
composition/decomposition strategy \emph{cannot manufacture a local \(\omega\)
below all of its inputs}: if every surrounding shape it can draw on has higher
\(\omega\), the arithmetic combination does too. Pushing a local \(\omega\)
\emph{below} its neighbourhood needs a mechanism that is not a blend of
neighbours --- one that rewrites the scheme's own product set rather than gluing
fixed sub-schemes together. Flip-graph and meta-flip-graph
moves~\cite{perminov2025fast,kauersmoosbauerwood2026} are exactly such a
mechanism, which is why they reach ranks our compositional closure cannot.

We must qualify this picture in one direction. The serendipitous
borrow-and-correct product (Section~\ref{ssec:serendipitous}), and more
speculatively projection (Section~\ref{ssec:decomposition}), are \emph{not}
purely additive blends: the borrow injects a correction term that lives outside
the sum of the consumed ranks, so they can in principle produce a local
\(\omega\) below all of their neighbours. What is unclear to us is
\emph{propagation}: even a genuinely lowered local \(\omega\) only helps the
broader landscape if the surrounding shapes in turn consume the improved scheme,
and we have no argument that a serendipitous local win propagates outward the way
a flip-graph improvement does. We therefore treat it as a local-improvement
mechanism of unproven global reach, not as a closure-completing one.

\subsection{Axis concatenation (block-additive composition)}
\label{ssec:concat}

If two schemes share two of the three axes, they can be concatenated
along the third: a \(\nmpshape{n}{m}{p_1}\) scheme of rank \(r_1\)
and a \(\nmpshape{n}{m}{p_2}\) scheme of rank \(r_2\) compose into a
\(\nmpshape{n}{m}{p_1 + p_2}\) scheme of rank \(r_1 + r_2\), simply
by computing the left and right halves of the result independently.
Analogous concatenations exist along the \(n\) and \(m\) axes. This
is sometimes referred to as ``plain split along a dimension'';
it should not be confused with the $\tau$-theorem disjoint sums
of Section \ref{ssec:disjoint-sum}, which are additive over the
tensor but not over a single axis.

The operator is trivial but pulls weight: at frontier-closure time
it is the cheapest way to cover ``rectangular'' shapes
(\(\nmpshape{n}{m}{p}\) with one large axis) from cubic catalog
entries.

\subsection{Kronecker product (multiplicative composition)}
\label{ssec:kronecker}

Given a scheme \(A\) at \(\nmpshape{n_1}{m_1}{p_1}\) of rank \(r_1\)
and a scheme \(B\) at \(\nmpshape{n_2}{m_2}{p_2}\) of rank \(r_2\),
the Kronecker product \(A \otimes B\) is a scheme at
\(\nmpshape{n_1 n_2}{m_1 m_2}{p_1 p_2}\) of rank \(r_1 r_2\). This
is the classical recursive substitution behind Strassen's \(\omega
\le \log_2 7\) bound: applying Strassen recursively to a power-of-two
\(\nmpshape{2^k}{2^k}{2^k}\) target gives \(7^k\) bilinear products.

The catalog exposes Kronecker as a composition operator and the
search enumerates 2-fold Kronecker chains over the pool.
Higher-arity chains are derivable but not currently enumerated
exhaustively -- they appear via repeated 2-fold composition.

\subsubsection{Serendipitous product (bud decomposition)}
\label{ssec:serendipitous}

The serendipitous product is a multiplicative composition that, unlike
the Kronecker product, can land \emph{below} the naïve rank product
\(r_1 r_2\). It was introduced by Smith~\cite{smith2002fast} and
generalised by Sedoglavic; we follow the bud-based formalisation of
Perminov~\cite{perminov2026serendipitous} (Definitions
\(2.9\)--\(2.12\) and \(\S 2.6\)).
The same fusion was independently formalised by Kauers, Moosbauer and
Wood~\cite{kauersmoosbauerwood2026} under the name ``divide less,
conquer more'': recursive calls that share an input, or whose output is
used in several positions, are merged into a single larger matrix
multiplication---our \(U\)- and \(W\)-buds, respectively---so that fewer
recursion levels are needed and the effective exponent drops below what
the tensor rank alone suggests (they report
\(\omega(\nmpshape{6}{6}{6}) \le 2.8019\)). Notably, their route to
\emph{obtaining} bud-rich decompositions is a \emph{search}: a flip-graph
search maximising the number of shareable \(\nmpshape{1}{1}{k}\),
\(\nmpshape{1}{\ell}{1}\), \(\nmpshape{m}{1}{1}\) sub-blocks under a
non-overlap constraint---the same objective as our bud-richness
selection---supplemented by \emph{random} elements of de~Groote's
symmetry group for support reduction; their \(\nmpshape{6}{6}{6}\)
scheme is obtained by analysing a Moosbauer--Poole rank-\(153\)
decomposition and re-searching for one richer in buds, and they report
no general \emph{construction} of such structure. Whether a deterministic
constructor (a targeted basis change in the spirit of the de~Groote
solve used for kin-row unification by Schwartz and
Zwecher~\cite{schwartzzwecher2025}) could replace this search is open.

The key object is a \emph{bud}: a pair (or group) of rank-one terms of
a scheme that share the \emph{same} \(u\) vector (or the same \(v\), or
the same \(w\)), where ``same'' is up to scaling, since a rank-one
tensor \(u \otimes v \otimes w\) is invariant under
\(u \mapsto \lambda u,\ w \mapsto \lambda^{-1} w\). Buds are
basis-free: the shared vector may be any linear combination of input
(or output) entries, not a single coordinate. Grouping the \(r\) terms
of a scheme by their buds decomposes it into \emph{elementary matmul
tensors}
\[
  \mathcal{T} \;=\; \sum_{i} S_i \, \nmpshape{N_i}{M_i}{P_i},
  \qquad
  r \;=\; \sum_{i} S_i \, N_i M_i P_i ,
\]
where a \(U\)-bud of \(k\) terms is an \(\nmpshape{1}{1}{k}\) block, a
\(V\)-bud a \(\nmpshape{k}{1}{1}\) block, a \(W\)-bud a
\(\nmpshape{1}{k}{1}\) block, combined buds give blocks such as
\(\nmpshape{2}{1}{2}\), and every term in no bud is a trivial
\(\nmpshape{1}{1}{1}\).

Given a second scheme \(\nmpshape{n_2}{m_2}{p_2}\), the serendipitous
product applies it to each elementary block independently:
\[
  \mathcal{T} \otimes \nmpshape{n_2}{m_2}{p_2}
  \;=\; \sum_{i} S_i \, \nmpshape{N_i n_2}{M_i m_2}{P_i p_2},
\]
and -- crucially -- each enlarged block is realised by the \emph{best
known} scheme for its format rather than by the naïve product, so
\[
  r_s \;=\; \sum_{i} S_i \cdot R(\nmpshape{N_i n_2}{M_i m_2}{P_i p_2}).
\]
A saving over the Kronecker product arises whenever
\(R(\nmpshape{N_i n_2}{M_i m_2}{P_i p_2}) < N_i M_i P_i \cdot
R(\nmpshape{n_2}{m_2}{p_2})\) for at least one block; for trivial
blocks the two coincide, so a scheme with no buds reduces exactly to
the Kronecker product.

\paragraph{Worked example: \(\nmpcount{}{4}{8}{12}{272}\).}
The Hopcroft--Kerr \(\nmpcount{}{2}{3}{4}{20}\) scheme has four
\(U\)-buds among its twenty terms, so it decomposes as
\(12\cdot\nmpshape{1}{1}{1} + 4\cdot\nmpshape{1}{1}{2}\) (note
\(12 + 4\cdot 2 = 20\)). Taking the serendipitous product with
\(\nmpshape{2}{4}{2}\) and using \(R(\nmpshape{2}{4}{2})=14\) and
\(R(\nmpshape{2}{4}{4})=26\),
\[
  \nmpshape{2}{3}{4}\otimes\nmpshape{2}{4}{2}
  \;=\; 12\cdot\nmpshape{2}{4}{2} + 4\cdot\nmpshape{2}{4}{4},
  \qquad
  r_s = 12\cdot 14 + 4\cdot 26 = 272,
\]
a scheme for \(\nmpshape{4}{12}{8}\) (equivalently
\(\nmpshape{4}{8}{12}\) up to axis permutation), below the Kronecker
value \(20\cdot 14 = 280\). The shorthand
``\(({\nmpshape{2}{3}{4}}{:}20 - 8)\otimes\nmpshape{2}{4}{2} +
4\,\nmpshape{2}{4}{4}\)'' used by the FMM catalog \cite{fmmlille}
records exactly this: eight of the twenty terms live in buds and are
promoted to the four \(\nmpshape{2}{4}{4}\) blocks, the other twelve
stay as \(\nmpshape{2}{4}{2}\).

\paragraph{Minimal example, multiplication by multiplication.}
The smallest serendipitous saving \emph{derives Strassen}. Take the base
\(\nmpshape{1}{1}{2}\) --- a row times a \(1\times 2\) row, products \(a\,b_1\)
and \(a\,b_2\); both share the single \(A\)-form \(a\), so they form one
\(U\)-bud of size \(2\). Multiplying by \(\nmpshape{2}{2}{1}\) gives target
\(\nmpshape{2}{2}{2}\), and the bud fuses the two inner copies into one
enlarged \(\nmpshape{2}{2}{2}\):
\[
  r_s = 1\cdot R(\nmpshape{2}{2}{2}) = 7 \;<\; 2\cdot R(\nmpshape{2}{2}{1}) = 8 .
\]
As a plain multiplication list, the Kronecker product runs the two
\(\nmpshape{2}{2}{1}\) matrix--vector blocks independently --- eight products,
the naïve \(\nmpshape{2}{2}{2}\) --- and both blocks reuse the \emph{same} \(A\)
entries; that reuse \emph{is} the bud:
\begin{quote}\small\ttfamily
p1=a11*b11\ p2=a12*b21\ p3=a21*b11\ p4=a22*b21\ \ (c11=p1+p2,\ c21=p3+p4)\\
p5=a11*b12\ p6=a12*b22\ p7=a21*b12\ p8=a22*b22\ \ (c12=p5+p6,\ c22=p7+p8)
\end{quote}
Serendipity replaces them with Strassen's seven, in which the shared \(A\) lets
three products span \emph{both} output columns at once:
\begin{quote}\small\ttfamily
M1=(a11+a22)(b11+b22)\ M2=(a21+a22)b11\ M3=a11(b12-b22)\\
M4=a22(b21-b11)\ M5=(a11+a12)b22\ M6=(a21-a11)(b11+b12)\\
M7=(a12-a22)(b21+b22)
\end{quote}
Exactly one product disappears (\(8\to 7\)): the bud's shared \(A\)-form is what
lets the cross-column combinations \(M_1,M_6,M_7\) absorb the fourth, otherwise
redundant pair. This is the whole mechanism in miniature --- buds expose
\emph{which} products can be merged, and a sub-additive enlarged block
(\(\nmpshape{2}{2}{2}\!:\!7<8\)) is what turns the merge into a saving.

\paragraph{Why bud-richness can be worth a higher rank.}
The rank-\emph{minimal} scheme for a format is often the \emph{wrong} base
for a serendipitous product. A plain example with catalogued ranks builds
\(\nmpshape{6}{6}{6}\) as \(\nmpshape{2}{2}{2}\otimes\nmpshape{3}{3}{3}\),
using \(R(\nmpshape{3}{3}{3})=23\) (Laderman) and \(R(\nmpshape{3}{3}{6})=40\)
(Smirnov; note \(40<2\cdot 23\)):
\begin{itemize}
\item The rank-minimal outer base, Strassen \(\nmpcount{}{2}{2}{2}{7}\), has
\emph{no} buds (its seven products use pairwise non-proportional
\(u,v,w\)), so it decomposes as \(7\cdot\nmpshape{1}{1}{1}\) and its
serendipitous product is just the Kronecker product,
\(7\cdot R(\nmpshape{3}{3}{3}) = 7\cdot 23 = 161\).
\item The higher-rank \emph{naïve} \(\nmpcount{}{2}{2}{2}{8}\) base is
bud-rich: each of its four \(A\)-entries feeds two outputs, giving four
\(U\)-buds, \(4\cdot\nmpshape{1}{1}{2}\). Its serendipitous product fuses
each bud with the inner \(\nmpshape{3}{3}{3}\) into a \(\nmpshape{3}{3}{6}\):
\(4\cdot R(\nmpshape{3}{3}{6}) = 4\cdot 40 = 160\).
\end{itemize}
The bud-rich rank-\(8\) base beats the rank-minimal rank-\(7\) base
(\(160<161\)) --- but \emph{only} through serendipity: its plain Kronecker
product costs \(8\cdot 23 = 184\), the worst of the three. A catalog that
keeps only the rank-minimal scheme per format discards exactly the base that
wins here; recovering the \(160\) requires storing the bud-rich rank-\(8\)
variant alongside the rank-\(7\) one.

\paragraph{How operations act on buds.}
Bud structure is not preserved uniformly by the composition operators, which
is why it is tracked explicitly rather than recomputed from the rank-minimal
scheme.
\begin{itemize}
\item \emph{Rank reduction destroys buds.} A bud exists because several
products share an input or output form; the standard route to lower rank
replaces un-mixed products by independent linear \emph{combinations} of
entries, making the \(u,v,w\) vectors pairwise non-proportional.
\(\nmpcount{}{2}{2}{2}{7}\) is the extreme case --- zero buds, versus four for
\(\nmpcount{}{2}{2}{2}{8}\); empirically every rank-minimal small-format
scheme we hold is bud-free, and buds appear only on larger, less-mixed
schemes.
\item \emph{The Kronecker product preserves and multiplies buds:} a term
\(u^1_i\otimes u^2_j\) is proportional to \(u^1_{i'}\otimes u^2_{j'}\) iff
both factors are proportional, so each bud of one factor lifts to a bud of
the product. The serendipitous product is the bud-aware way of taking this
Kronecker product.
\item \emph{Axis flip / permutation (the \(S_3\) orbit) permutes the bud
type:} a \(U\)-bud becomes a \(V\)- or \(W\)-bud under the corresponding
relabelling, leaving the bud \emph{count} unchanged --- one orbit already
exposes bud-richness on all three axes.
\item \emph{Projection / DCE can only shrink buds} (it deletes terms), while
concatenating a scheme with a copy of itself \emph{creates} \(U\)-buds (each
product gains a twin with the same \(u\) but a fresh output column).
\end{itemize}

\paragraph{Which schemes to retain.}
Because \(r_s\) depends only on the bud \emph{multiset}
\(\{S_i\,\nmpshape{N_i}{M_i}{P_i}\}\), two schemes with the same multiset are
interchangeable for every second factor: the store keys on the multiset
(with rank and additions as tie-breakers), not on scheme identity. By the
\(S_3\) isotropy above, a ``best-\(U\), best-\(V\), best-\(W\)'' triple is
redundant --- one orbit representative covers all three axes, with the
orientation chosen at product time to align the bud axis with the contracted
axis of the second factor. What must be retained is the \emph{Pareto
frontier} in \((\text{rank},\,\text{bud multiset})\): a higher-rank scheme
is worth keeping iff its richer buds win some serendipitous product the
lower-rank scheme cannot, while a scheme of no-better rank whose bud
multiset is contained in another's is dominated and may be dropped.
Independent schemes sharing the same maximal bud multiset collapse to the
cheapest representative.

\paragraph{Aligning buds under isotropy.}
Because the Kronecker product multiplies proportionalities, the bud structure
of \(\mathcal{T}_1\otimes\mathcal{T}_2\) depends on how the two factors' bud
axes line up. A \(U\)-bud of size \(a\) in \(\mathcal{T}_1\) \emph{aligned}
with a \(U\)-bud of size \(b\) in \(\mathcal{T}_2\) (same axis) yields a
\(U\)-bud of size \(a\,b\) in the product, since a shared product \(u\)
requires both \(u^1_i\propto u^1_{i'}\) and \(u^2_j\propto u^2_{j'}\); a
crossed pair (\(U\) against \(V\)) reinforces nothing and leaves only the
size-\(a\) bud. The only isotropy that can realign buds is the discrete
\(S_3\) axis permutation: a basis change \(u\mapsto M u\) preserves
proportionality (\(M u_i\propto M u_j \Leftrightarrow u_i\propto u_j\)), so
the continuous \(GL\) part of the isotropy group leaves the bud partition
\emph{invariant} and only the axis relabelling moves buds between \(U\),
\(V\) and \(W\). Two levers follow. For the (asymmetric) serendipitous
product \(\mathcal{T}_1\otimes\nmpshape{n_2}{m_2}{p_2}\), where only
\(\mathcal{T}_1\)'s buds are used, orienting \(\mathcal{T}_1\) by \(S_3\)
selects which target axis carries the \(k\)-fold enlargement --- a \(U\)-bud
gives \(\nmpshape{n_2}{m_2}{k p_2}\), a \(V\)-bud
\(\nmpshape{k n_2}{m_2}{p_2}\), whose \(R\) differ --- and so changes
\(r_s\) directly; the search should try all three orientations of the
bud-provider. For a full \(\mathcal{T}_1\otimes\mathcal{T}_2\) reused as a
base, aligning both factors' bud axes maximises the \(a\,b\) reinforcement
and yields a bud-richer scheme for the next multiplication.

\paragraph{What it is, and is not.}
The serendipitous product is \emph{exact} and \emph{combinatorial}: it
reads buds off the base scheme's factor matrices and looks up
sub-block ranks; it requires neither numerical optimisation nor border
rank. It differs from disjoint-sum / \(\tau\) decomposition
(\S\ref{ssec:disjoint-sum}), which splits the \emph{target} tensor,
and from pair fusion (\S\ref{ssec:pair-fusion}), which fuses two
equal-shape sub-products: here a single base scheme is multiplied by a
second, and the saving comes from grouping the base's own terms into
larger matmul blocks before multiplying. Our bud recognizer and
bud-block constructor follow Perminov's \(\S 2.6\) construction; the
assembled scheme is certified by the exact symbolic verifier.

\paragraph{Smoothness governs the serendipitous hit rate.}
A serendipitous product needs the target to factor as
\(\nmpshape{n_1}{m_1}{p_1}\otimes\nmpshape{n_2}{m_2}{p_2}\) with \emph{both}
factors non-trivial, and the number of such factorisations of
\(\nmpshape{N}{M}{P}\) is \(d(N)\,d(M)\,d(P)\), the product of the per-axis
divisor counts. A \emph{smooth} format --- one whose sides have many (distinct)
prime factors, e.g. \(30=2\cdot3\cdot5\) --- therefore exposes far more
base/inner splits than a prime or prime-power format, and each split is an
independent chance for a bud-rich base to meet a sub-additive enlarged inner
(both conditions must hold at once). One expects, and we observe, that
serendipity wins disproportionately on smooth formats. Writing \(\Omega(NMP)\)
for the number of prime factors of \(N\,M\,P\) counted with multiplicity,
Table~\ref{tab:serendip-smooth} reports, over all \(5456\) non-commutative
formats in the catalogue, how often the \emph{best} scheme we hold is the
serendipitous one.

\begin{table}[t]
\centering\small
\begin{tabular}{rrrr}
\toprule
\(\Omega(NMP)\) & formats & serendipity best & \% \\
\midrule
\(\le 5\) & 2013 & 0 & 0.0 \\
6 & 1122 & 6 & 0.5 \\
7 & 979 & 12 & 1.2 \\
8 & 654 & 11 & 1.7 \\
9 & 389 & 9 & 2.3 \\
10 & 178 & 3 & 1.7 \\
\(\ge 11\) & 121 & 0 & 0.0 \\
\midrule
all & 5456 & 41 & 0.75 \\
\bottomrule
\end{tabular}
\caption{How often the best known scheme for a format is a serendipitous
product, bucketed by the smoothness \(\Omega(NMP)\) of the format. No
serendipitous scheme wins below \(\Omega=6\); above that floor the hit rate
rises with smoothness and peaks near \(\Omega=9\). Formats where serendipity
wins have mean \(\Omega = 7.78\), against \(6.29\) for the whole population.}
\label{tab:serendip-smooth}
\end{table}

Below \(\Omega=6\) no serendipitous scheme is ever best: the smallest
non-trivial base/inner pair already needs \(\Omega\ge6\) (e.g.
\(\nmpshape{2}{3}{7}\otimes\nmpshape{3}{3}{3}\)), so smaller formats admit no
serendipitous decomposition at all. Above the floor the win rate tracks
smoothness, peaking near \(\Omega=9\). The fall-off at \(\Omega\ge11\) is a
coverage artefact --- those formats are large, beyond our materialisation cap,
where fewer serendipitous schemes are yet built --- not a ceiling on the
method. This concerns the practical per-format implied exponent, not the
asymptotic \(\omega\); operationally it gives the search a prioritisation rule:
spend serendipitous effort on smooth formats first, and treat prime cubes as
direct-construction targets where the method cannot help.

\subsection{Decomposition: block-splitting and the sub-problem multiset}
\label{ssec:decomposition}

The operators above build \emph{up}; the rest of this section builds
\emph{down}. \emph{Decomposition} splits a large \(\nmpshape{N}{M}{P}\)
product into a grid of blocks and computes that grid with a smaller
\emph{outer} scheme of shape \(\nmpshape{n_o}{m_o}{p_o}\) and rank
\(r_o\): cut each target axis into that many parts, view the target as an
\(n_o\times m_o\) by \(m_o\times p_o\) block product, and let each of the
outer scheme's \(r_o\) bilinear products compute one block product. The
basic expectation to hold onto: a \(\nmpshape{2}{2}{2}\) outer scheme
turns the target into \emph{seven} block products --- seven
sub-problems --- and Laderman's \(\nmpshape{3}{3}{3}\) into twenty-three.

\paragraph{The sub-problem multiset.} List those sub-problems by shape,
with multiplicity, and you get the \emph{sub-problem multiset} of the
decomposition. Since each is computed independently, the realised rank is
just the sum of their ranks,
\[
  R(\nmpshape{N}{M}{P}) \;\le\;
  \sum_{\sigma\in\text{multiset}} R(\sigma),
\]
so the multiset is the \emph{only} thing that fixes the cost --- order
and labelling are irrelevant. When the parts are all equal the seven (or
twenty-three) sub-problems are identical, the multiset is one shape
repeated, and the sum collapses to \(r_o\cdot R(\text{block})\) ---
exactly the Kronecker product (Section~\ref{ssec:kronecker}). The
interesting regime is \emph{unbalanced} parts, where the sub-problems
differ. (This is a different multiset from the \emph{bud} multiset of
Section~\ref{ssec:serendipitous}, which groups a scheme's own rank-one
terms; both are ``multisets'' but of different things, and a more
specific name for one of them may be worth adopting.)

\paragraph{Padding and the max.} The realised size of each sub-product is
where padding enters, and it is the heart of the matter. When a product
\emph{sums} two blocks of unequal size on an output axis, the smaller is
embedded into the larger with zeros, so it is realised at the \emph{max}
of the two parts on that axis (the contraction axis takes the
\emph{min}); the padding zeros are never multiplied. A product that sums
across an axis is therefore charged that axis's \emph{large} part, while a
product that uses a \emph{single} block stays at the part it touches.
Hence the expensive corner \(\nmpshape{M_1}{N_1}{P_1}\) is paid by every
all-sum product, and the cheap all-small corner
\(\nmpshape{M_2}{N_2}{P_2}\) is reachable only by a product that is
single-block on \emph{every} axis.

It is this max that decides Strassen versus Winograd. At a fixed
orientation both induce the \emph{same} multiset (Section
\ref{ssec:strassen-vs-winograd}: the \(\nmpshape{3}{3}{3}\) example gives
\(30\) either way) and differ only in additions. But Strassen's seven
products are all two-block sums, so every one pays the max on every axis
and the all-small corner never appears; Winograd has two \emph{single}-block
products, which the axis-flip orbit (Section~\ref{ssec:axis-flip}) can
steer onto the small corner, removing one max. So the padding-max --- not
the rank, not the additions --- is \emph{why} Winograd beats Strassen on
unbalanced splits. \textbf{Which template} and \textbf{which orientation}
are the two knobs; recombination with allocation
(Section~\ref{ssec:recombination}) searches over them.

\subsubsection{Strassen, Winograd, and the structural difference}
\label{ssec:strassen-vs-winograd}

The two canonical rank-7 algorithms for \(\nmpfield{\Reals}{2}{2}{2}\)
are Strassen 1969 and Winograd's 1971 reformulation. This section gives
the \emph{generic} effect of using either as a \(\nmpshape{2}{2}{2}\)
template over a split
\(\nmpshape{M_1{+}M_2}{N_1{+}N_2}{P_1{+}P_2}\): the seven sub-products it
induces, sized by zero-sparing. At this simple (fixed-orientation) level
the two schemes induce the \emph{same} multiset of sub-problems and
differ only in additions (Strassen 1969: 18; Winograd 1971: 15). The
\emph{axis-flip orbit} of Section~\ref{ssec:axis-flip} then re-orients
the template -- which \emph{does} change the multiplications, and is
where the two schemes diverge.

\paragraph{Heterogeneous split and zero-sparing.}
Project a target onto a \(\nmpshape{2}{2}{2}\) template by splitting
each axis in two, \(M = M_1 + M_2\), \(N = N_1 + N_2\),
\(P = P_1 + P_2\), with the larger part first and the smaller part at
the bottom/right (real data top-left, any padding at the bottom). The
four \(A\)-blocks then have shapes \(A_{11}\!:\!M_1\times N_1\),
\(A_{12}\!:\!M_1\times N_2\), \(A_{21}\!:\!M_2\times N_1\),
\(A_{22}\!:\!M_2\times N_2\), and likewise for \(B\) and \(C\). Each
product is a (sum of \(A\)-blocks)\(\,\cdot\,\)(sum of \(B\)-blocks); to
add blocks of unequal shape the smaller is embedded into the larger
envelope with zeros. \emph{Those zeros are never multiplied}: a product
is realised at the size of its genuine nonzero overlap -- rows = the max
real rows of the left sum, columns = the max real columns of the right
sum, inner = the \emph{min} of the two sides. A product that touches
only the smaller part of an axis is thereby reduced to a lower rank.

\paragraph{The seven products, sized (canonical labelling).}
With block ``11'' the larger part on every axis, each product's realised
shape under zero-sparing (``blk'' = single block, ``sum'' = two-block
sum) is:

\begin{center}\small
\begin{tabular}{@{}llll@{}}
\toprule
\multicolumn{2}{@{}l}{\textbf{Strassen}} & operands & realised size \\
\midrule
\(M_1\) & \((A_{11}{+}A_{22})(B_{11}{+}B_{22})\) & sum\(\cdot\)sum & \(\nmpshape{M_1}{N_1}{P_1}\) \\
\(M_2\) & \((A_{21}{+}A_{22})\,B_{11}\)          & sum\(\cdot\)blk & \(\nmpshape{M_2}{N_1}{P_1}\) \\
\(M_3\) & \(A_{11}(B_{12}{-}B_{22})\)            & blk\(\cdot\)sum & \(\nmpshape{M_1}{N_1}{P_2}\) \\
\(M_4\) & \(A_{22}(B_{21}{-}B_{11})\)            & blk\(\cdot\)sum & \(\nmpshape{M_2}{N_2}{P_1}\) \\
\(M_5\) & \((A_{11}{+}A_{12})\,B_{22}\)          & sum\(\cdot\)blk & \(\nmpshape{M_1}{N_2}{P_2}\) \\
\(M_6\) & \((A_{21}{-}A_{11})(B_{11}{+}B_{12})\) & sum\(\cdot\)sum & \(\nmpshape{M_1}{N_1}{P_1}\) \\
\(M_7\) & \((A_{12}{-}A_{22})(B_{21}{+}B_{22})\) & sum\(\cdot\)sum & \(\nmpshape{M_1}{N_2}{P_1}\) \\
\bottomrule
\end{tabular}
\end{center}

\begin{center}\small
\begin{tabular}{@{}llll@{}}
\toprule
\multicolumn{2}{@{}l}{\textbf{Winograd}} & operands & realised size \\
\midrule
\(P_1\) & \(A_{11}\,B_{11}\)                     & blk\(\cdot\)blk & \(\nmpshape{M_1}{N_1}{P_1}\) \\
\(P_2\) & \(A_{12}\,B_{21}\)                     & blk\(\cdot\)blk & \(\nmpshape{M_1}{N_2}{P_1}\) \\
\(P_3\) & \((A_{21}{+}A_{22})(B_{12}{-}B_{11})\) & sum\(\cdot\)sum & \(\nmpshape{M_2}{N_1}{P_1}\) \\
\(P_4\) & \(S_2\,T_2\)                           & sum\(\cdot\)sum & \(\nmpshape{M_1}{N_1}{P_1}\) \\
\(P_5\) & \(S_3\,T_3\)                           & sum\(\cdot\)sum & \(\nmpshape{M_1}{N_1}{P_2}\) \\
\(P_6\) & \(S_4\,B_{22}\)                        & sum\(\cdot\)blk & \(\nmpshape{M_1}{N_2}{P_2}\) \\
\(P_7\) & \(A_{22}\,T_4\)                        & blk\(\cdot\)sum & \(\nmpshape{M_2}{N_2}{P_1}\) \\
\bottomrule
\end{tabular}
\end{center}
\noindent with \(S_2{=}A_{21}{+}A_{22}{-}A_{11}\),
\(S_3{=}A_{11}{-}A_{21}\),
\(S_4{=}A_{11}{+}A_{12}{-}A_{21}{-}A_{22}\),
\(T_2{=}B_{11}{-}B_{12}{+}B_{22}\), \(T_3{=}B_{22}{-}B_{12}\),
\(T_4{=}B_{11}{-}B_{12}{-}B_{21}{+}B_{22}\).

\paragraph{The generic multiset.}
Collecting either ``realised size'' column gives the same multiset for
both schemes:
\[
  2\,\nmpshape{M_1}{N_1}{P_1}
  + \nmpshape{M_2}{N_1}{P_1}
  + \nmpshape{M_1}{N_1}{P_2}
  + \nmpshape{M_1}{N_2}{P_1}
  + \nmpshape{M_1}{N_2}{P_2}
  + \nmpshape{M_2}{N_2}{P_1}.
\]
The largest corner \(\nmpshape{M_1}{N_1}{P_1}\) appears twice, the
all-small corner \(\nmpshape{M_2}{N_2}{P_2}\) not at all, and the split
is \emph{not} symmetric: the contraction axis \(N\) carries its smaller
part in three of the seven products while \(M\) and \(P\) carry theirs in
two each -- so one minimises the small-part products by making \(N\) the
largest axis. With equal parts (even dimensions) the multiset collapses
to \(7\,\nmpshape{M/2}{N/2}{P/2}\), the textbook recursion. At this fixed
orientation, then, the choice of Strassen vs Winograd changes only the
additive schedule, not the multiplications; re-orienting the template
(Section~\ref{ssec:axis-flip}) is what changes the multiplications.

\paragraph{Worked example: \(\nmpshape{3}{3}{3}\).}
Split \(3 = 2{+}1\) on every axis (\(M_1{=}N_1{=}P_1{=}2\),
\(M_2{=}N_2{=}P_2{=}1\)). Both schemes yield
\[
  2\,\nmpshape{2}{2}{2} +
  \nmpshape{1}{2}{2} + \nmpshape{2}{2}{1} + \nmpshape{2}{1}{2} +
  \nmpshape{2}{1}{1} + \nmpshape{1}{1}{2},
\]
i.e.\ \(2\cdot 7 + 4 + 4 + 4 + 2 + 2 = 30\) leaf multiplications
(recursing each \(\nmpshape{2}{2}{2}\) with the same rank-7 scheme) --
against \(7\cdot 7 = 49\) for the naive pad-to-even \(\nmpshape{4}{4}{4}\)
embedding and \(3^3 = 27\) for the schoolbook product. The
\(49 \to 30\) gap is the zero-sparing; the further drop available by
re-orienting the template is deferred to Section~\ref{ssec:axis-flip}.

\paragraph{Why one needs a \emph{set} of bases.} Strassen and Winograd
are only two points in the space of rank-7 \(\nmpshape{2}{2}{2}\) schemes,
and --- as the axis-flip orbit will show --- what matters at an unbalanced
split is the per-product block/sum structure (how many single-block
products a scheme has, hence which sub-problem multisets its orbit can
reach). To exploit this fully the pool should carry a \emph{set} of
rank-7 \(\nmpshape{2}{2}{2}\) bases that are inequivalent in their
multiset behaviour, not just the two textbook ones; generating that set
means enumerating the rank-7 normal forms (Heun 1994; de~Groote's
classification) and keeping one representative per axis-flip orbit
(Section~\ref{ssec:axis-flip}). For \(\nmpshape{2}{2}{2}\) this enumeration
is small and we do it exhaustively. The same idea applies to any base
shape --- a richer base set gives the recombination more multisets to
choose from --- but the space of rank-\(r\) schemes at larger shapes
(\(\nmpshape{3}{3}{3}\) and up) is far too large to enumerate, so there we
fall back on the imported/known schemes plus the axis-flip orbit of each.

\subsubsection{Axis-flip orbit}
\label{ssec:axis-flip}

For any bilinear scheme \((U, V, W)\) of shape
\(\nmpshape{n}{m}{p}\), reversing the row order of any of the three
factor matrices gives another valid scheme of the same shape and the
same rank. Composing the three independent axis-flips yields an
orbit of size 8 per scheme (the elementary abelian \(2^3\)).

Pure axis-flip is a trivial relabeling under uniform allocations
(the orbit collapses) but becomes non-trivial under \emph{unbalanced}
allocations such as the \((9, 8) \times \ldots\) split above.
Different axis-flip orbit members give rise to different padding
patterns and hence different total ranks. Section
\ref{ssec:recombination} explains why this matters at the
recombination step: an axis-flip of the outer scheme corresponds to
permuting which block-index in the target receives the padding
penalty.

We canonicalise the pool of outer schemes by axis-flip orbit
(picking one representative per orbit) and then re-enumerate the
mask at search time, so the on-disk catalog stays compact while the
search retains full coverage. The mask enumeration uses an
analytical short-cut (per-product block-support bitmasks; see
Section \ref{sec:search}) so 8-mask coverage costs only a constant
factor over single-mask evaluation.

\paragraph{Why the orbit separates Strassen from Winograd.}
This is exactly where the two rank-7 templates of
Section~\ref{ssec:strassen-vs-winograd} stop being interchangeable. At a
fixed orientation they share the generic multiset above -- two full
\(\nmpshape{M_1}{N_1}{P_1}\) products and no all-small
\(\nmpshape{M_2}{N_2}{P_2}\). Strassen's seven products are all
two-block sums, so under \emph{any} orbit member each still spans a full
part on every axis: Strassen keeps its two full products and can never
realise the all-small corner. Winograd's two single-block products
(\(A_{11}B_{11}\), \(A_{12}B_{21}\)) \emph{can} be flipped onto the small
blocks -- one becomes \(A_{22}B_{22} = \nmpshape{M_2}{N_2}{P_2}\) -- so
the best Winograd orbit member drops to a single full product and a
strictly smaller multiset whenever the split is unbalanced (the orbit
collapses, and the gain vanishes, only when the parts are equal).
Concretely, \(\nmpshape{17}{17}{17}\) at the \((9,8)^3\) allocation costs
\(2930\) multiplications for the best Winograd orbit member -- its lone
all-small \(\nmpshape{8}{8}{8}\) product replacing one
\(\nmpshape{9}{9}{9}\) -- against \(2940\) for Strassen. It is a genuine
\emph{multiplication} saving, distinct from Winograd's lower addition
count.

\subsection{Recombination with allocation}
\label{ssec:recombination}

The central composition operator in the catalog. It uses a
\emph{small} catalog scheme --- the outer \emph{atom}, e.g.\ Strassen
\(\nmpshape{2}{2}{2}\) --- to multiply a \emph{larger} target
\(\nmpshape{n}{m}{p}\) by block-partitioning the target onto the
atom's grid. The atom's dimensions \(\nmpshape{n_o}{m_o}{p_o}\) are
small (typically \(2\) or \(3\), and \(\le n, m, p\)) and fix only
\emph{how many} blocks each target axis is cut into, not their sizes.
Concretely, given the atom of rank \(r_o\), choose three
\emph{allocations}: a partition of the target's \(n\) into \(n_o\)
parts, of \(m\) into \(m_o\) parts, and of \(p\) into \(p_o\) parts.
The atom then treats the target as an \(n_o\times m_o\) by
\(m_o\times p_o\) block-matrix product: each of its \(r_o\) outer
products becomes a smaller matmul on the corresponding blocks --- a
sub-shape fixed by the allocation --- which is looked up in the
catalog as a sub-scheme, and the assembled algorithm has rank
\(\sum_k r_{\text{sub-}k}\).

Uniform allocation reduces to Kronecker product (Section
\ref{ssec:kronecker}); non-uniform allocation gives strictly more
flexibility -- e.g.\ targeting \(\nmpshape{17}{17}{17}\) via
Strassen with the allocation \((9, 8)^3\). Non-uniform recombination
covers the gap left by Kronecker chains, at the price of
\emph{padding} -- a sub-shape derived from two unequal blocks
inherits the size of the larger block on the relevant axis.

\subsection{Projection margin: a death-rate score}
\label{ssec:projmargin}

Projection is the downward operator: restrict a larger scheme to a smaller
shape and dead-code-eliminate the products that die. A higher-rank parent
can thereby project to a \emph{lower}-rank child, and this has a clean
closed form worth isolating, because it turns ``keep the structurally
better scheme, not the lower-rank one'' into an explicit, computable
criterion --- the downward analogue of the bud score for the serendipitous
product (Section~\ref{ssec:serendipitous}).

A product survives the restriction iff \emph{each} of its three factors
\(U_k,V_k,W_k\) still has a nonzero inside the kept sub-block; it is dead-code-
eliminated as soon as \emph{any one} of them vanishes there. Fix an axis, say
the first (\(n\)); only the two matrices carrying an \(n\)-index matter --- \(U\)
(input rows \(i\)) and \(W\) (output rows \(i\)). Write \(U_n(k)\) and \(W_n(k)\)
for the sets of \(n\)-indices column \(k\) of \(U\), resp.\ \(W\), touches.
Dropping index \(i\) kills product \(k\) exactly when \emph{either} its
\(\alpha\)-form collapses (\(U_n(k)=\{i\}\): all of \(U_k\) sits in the dropped
input row, so \(\alpha_k\equiv0\) once that row is zeroed) \emph{or} it writes
only the dropped output (\(W_n(k)=\{i\}\)). The \emph{private rank} at \(i\) is
the count of such products,
\[
  \rho_n(i)=\#\bigl\{\,k:\ U_n(k)=\{i\}\ \text{ or }\ W_n(k)=\{i\}\,\bigr\},
\]
(an \(\text{or}\), not an \(\text{and}\): a product with \(U_n(k)=\{i\}\),
\(W_n(k)=\{i'\}\) dies under dropping \(i\) \emph{or} \(i'\)). Each axis uses two
of the three matrices --- \(n\!\to\!\{U,W\}\), \(m\!\to\!\{U,V\}\),
\(p\!\to\!\{V,W\}\) --- and the \emph{projection margin} of \(S\) (single index,
best axis) is
\[
  \mu(S)=\max_{A\in\{n,m,p\}}\ \max_{i}\ \rho_A(i).
\]
Projecting out the best single index then gives \emph{exactly}
\[
  R_{\text{after}} \;=\; R_{\text{before}}-\mu(S),
\]
and dropping a set \(D\) of \(\delta\) indices on one axis generalises to
\(\mu^{(\delta)}_A(S)=\max_{|D|=\delta}\#\{k:\mathrm{supp}_A(k)\subseteq D\}\),
with \(R_{\text{after}}=R_{\text{before}}-\mu^{(\delta)}\).

Projection therefore minimises \(R-\mu\), \emph{not} \(R\): a scheme is worth
keeping over a lower-rank rival whenever \(R_S-\mu(S)<R_{S'}-\mu(S')\), i.e.\ on
the \((\text{rank},\,\text{projection margin})\) Pareto frontier. The canonical
instance is \(\nmpshape{29}{30}{30}\): the rank-best parent, a flat
\(\nmpcount{}{30}{30}{30}{12688}\) (Schwartz--Zwecher 2025, an ALS-style
decomposition), has \(\mu=62\) and projects to \(12626\); the higher-rank
\emph{structured} \(\nmpcount{}{30}{30}{30}{12710}\) (Pan trilinear
aggregation) has \(\mu=122\) and projects to \(12588\), matching FMM. Structure
localises whole product-slabs to boundary slices (large \(\rho\)); a flat
decomposition spreads every product across all indices (small \(\rho\)), so it
projects worse despite the lower rank.

\paragraph{A small worked instance, and how to \emph{raise} the margin.}
Ordering our small catalog by \(\mu\), the high-margin schemes are uniformly the
\emph{structured} ones (Pan trilinear aggregation) and the \emph{flip-graph}
ones (Perminov), never the flat rank-minimal decompositions. For example the
Perminov \(\nmpcount{}{6}{6}{7}{183}\) has \(\mu=26\) on the \(p\)-axis and
projects to \(\nmpshape{6}{6}{6}\) at \(183-26=157\) --- acceptable, though not
the catalog optimum \(153\); the rank-minimal \(\nmpshape{6}{6}{7}\) scheme
projects worse. Two caveats sharpen the picture. First, at \emph{small}
dimension the flatter low-rank cube still wins projection outright, because the
rank gap dwarfs the margin gap; the structured-cube advantage only \emph{flips
the winner} once the two ranks converge, as at \(\nmpshape{30}{30}{30}\)
(\(12688\) vs \(12710\), decided by \(\mu=62\) vs \(122\)). Second --- and
parallel to the open question for buds --- it is not obvious how to
\emph{construct} a scheme of the \emph{same} rank but a higher \(\mu\). A
promising route is the flip graph: a meta-flip (a rank-preserving move in the
flip-graph of decompositions) leaves \(R\) fixed while reshuffling the product
supports, and can concentrate more products onto a single boundary index ---
i.e.\ search the flip-graph for the \(\mu\)-maximal representative of a rank
class, exactly as one would search it for a bud-maximal one. We leave this
generative direction to future work; here \(\mu\) is used to \emph{select}
among existing schemes, not to synthesise new ones.

\paragraph{Downward propagation is acyclic and cheap.} Unlike the serendipitous
product, whose candidate bases form a multiplicative (divisor) lattice,
projection's dependency graph is a DAG ordered by strictly decreasing dimension:
each scheme \(\nmpshape{N}{M}{P}\) projects only to its few \(\delta\)-neighbours
\(\nmpshape{N-\delta}{M}{P},\dots\). The catalog therefore need not be re-swept
to a fixed point round by round; a single newly-arrived scheme (an import, or a
fresh search win) with \(\mu>0\) triggers a bounded \emph{downward wave} ---
project it into its neighbours, and any neighbour that improves is itself
enqueued for its own projections. Termination is immediate (dimension strictly
decreases) and the work is event-driven: a new \(\nmpshape{n}{n}{n}\) cube
refreshes exactly the rectangular and odd-cube shapes beneath it, nothing more.
This makes projection the natural operator to fire on every catalog update
(e.g.\ each daily Perminov / FMM sync), in contrast to the global
recombination / serendipity passes.

\subsection{Zero-peeling: input- and output-side (DIS09 \(\gamma_5\))}
\label{ssec:peel}

When a recombination over-allocates an axis and pads the excess with
zeros, the padded rows/columns are dead and the affected sub-products
can be computed at a smaller shape. There are two directions, dual under
the matmul tensor's \(S_3\) symmetry:
\begin{itemize}[itemsep=2pt]
  \item \emph{input-side}: a sub-product reads \(A\) only in some rows, or
    \(B\) only in some columns (the rest are padding zeros) -- shrink its
    input extents;
  \item \emph{output-side}: a sub-product's \(W\)-support lands only in
    output positions that are peeled away after un-padding -- shrink its
    output extents. This is the Drevet--Islam--Schost 2009 / Islam 2009
    \(\gamma_5\) reduction.
\end{itemize}
The peel-aware recombination shrinks each sub-product to the
\emph{elementwise minimum} of both directions, and the block-split search
enumerates the (allocation, peel) patterns automatically (overshoot
bounded per axis). The canonical \(\gamma_5\) case -- \(\nmpshape{3}{3}{3}\)
via padded \(\nmpshape{4}{4}{4}\) Strassen, where the corner product
collapses \(\nmpshape{2}{2}{2}{=}7 \to \nmpshape{1}{2}{1}{=}2\) -- is
closed and tested.

Peeling is the second padding-fighter after axis-flip choice
(Section~\ref{ssec:axis-flip}); the two are complementary -- flip
distributes padding across products, peel removes it from a product whose
slice was dead anyway -- and, like axis-flip, it is \emph{not} a
sharing-discovery mechanism (Section~\ref{sec:nonoverlap}).

\paragraph{Limit.} Peeling alone does not close the odd-cubic gaps such
as \(\nmpshape{17}{17}{17}{=}2934\): there, no standard Strassen product
has its support inside the peeled corner, so peeling leaves us at
\(2940\). FMM-Lille's \(2934\) recipe instead pair-fuses the diagonal
\(\nmpshape{9}{9}{9}\) products and substitutes the corner with
\(\nmpshape{8}{8}{8}\) -- a pair-fusion (Section~\ref{ssec:pair-fusion})
plus outer-template substitution, not a peel.

\subsection{Pair fusion (Pan trilinear aggregation)}
\label{ssec:pair-fusion}

Pan 1980 \cite{pan1980} observed that two same-shape sub-products
arising at distinct positions in a larger scheme can be computed
\emph{jointly} via a trilinear aggregation that takes
\(a b c + a b + b c + c a\) bilinear products in place of the
naïve \(2 a b c\). Cyclic pair-fusion identifies pairs of outer
products with a compatible block-position pattern and substitutes a
fused sub-scheme. This is the first derivation strategy that
\emph{does} introduce sharing across outer products -- but only in a
constrained, pattern-matching way, not as a generic discovery
mechanism. The pair-fusion operator is implemented in
\code{PairFusedRecombination}.

\paragraph{When it pays: a rank-density / \(\omega\) threshold.}
Fusing is worthwhile only against the \emph{best} independent
computation, i.e. \(abc + ab + bc + ca < 2\,R(\nmpshape{a}{b}{c})\) -- not
against the naïve \(2abc\). Writing the leaf's \emph{rank density}
\(\rho = R(\nmpshape{a}{b}{c})/(abc)\), this is
\[
  \rho \;>\; \tfrac{1}{2}\Big(1 + \tfrac{ab+bc+ca}{abc}\Big)
  \;\approx\; \tfrac{1}{2}.
\]
So pair-fusion helps exactly when the leaf is still \emph{``fat''} -- its
rank has not been driven below about half the naïve \(abc\). For a cubic
leaf \(\nmpshape{k}{k}{k}=R\) with implied exponent
\(\omega_{\mathrm{leaf}} = \log_k R\) (naïve \(R=k^3\Rightarrow\omega=3\)),
the same condition reads
\[
  \omega_{\mathrm{leaf}} \;>\; \omega^\star
  := \log_k\!\tfrac{k^3+3k^2}{2}
  \;=\; 3 - \log_k 2 + \log_k\!\big(1+\tfrac{3}{k}\big)
  \;\approx\; 3 - \log_k 2 .
\]

\emph{Worked examples.} \(\nmpshape{7}{7}{7}=249\) has
\(\rho=249/343=0.726\) and \(\omega_{\mathrm{leaf}}=\log_7 249=2.835\),
above its threshold \(\omega^\star=\log_7 245=2.827\): it \emph{fuses}
(pair cost \(490 < 2\cdot 249 = 498\), saving \(8\)). By contrast
\(\nmpshape{8}{9}{9}=430\) has \(\rho=430/648=0.664\), just below its
threshold \(0.673\): it is \emph{not} fused (pair cost
\(873 > 2\cdot 430 = 860\)).

This is why pair-fusion is a \emph{small-base, constant-factor} device:
as \(k\to\infty\) the threshold \(\omega^\star\to 3\) while good leaves
have \(\omega_{\mathrm{leaf}}\) far below \(3\), so the saving vanishes
asymptotically. (The \emph{full} Pan trilinear aggregation -- aggregating
many products, not just pairs -- is the genuinely \(\omega\)-improving
technique; what the catalog fuses is only the two-product instance.)

\paragraph{Why pairs, not triples.}
Given several same-shape leaves, why fuse them two at a time rather than
three? Because there is \emph{no} ``fuse three \emph{separate} products''
primitive. The pair identity genuinely fuses two cyclically-related
products \(\nmpshape{a}{b}{c}+\nmpshape{b}{c}{a}\) at cost
\(abc+ab+bc+ca\). The \(\tfrac{1}{3}\) of the full Pan/Islam aggregation
does \emph{not} come from fusing three products -- it comes from one
product's own 3-fold cyclic index symmetry. A single
\(\nmpshape{n}{n}{n}\) costs \(\approx n^3/3 + O(n^2)\), and the
\(O(n^2)\) term is itself a \emph{best-of} several Pan-family closed
forms: \((n^3+12n^2+11n)/3\) (Islam 2009~\cite{islam2009} Prop.~1,
republished as DIS09~\S3) wins for moderate even \(n\), but the smaller
\(\tfrac{15}{4}n^2\) quadratic of the \(\tfrac{45}{4}n^2\) branch
(Hadas--Schwartz 1982, sharpened by Schwartz--Zwecher
2025~\cite{schwartzzwecher2025}) overtakes it for even \(n \gtrsim 28\);
odd \(n\) uses \((n^3+15n^2+14n-6)/3\). The chooser takes the
\emph{minimum} over all of these (\code{PanTrilinearAggregation.bestPanTaBound}),
so the large-\(n\) \(\tfrac{45}{4}n^2\) regime is accounted for, not just
the \(12n^2\) branch. So the real per-leaf choice is shallow pair fusion
(\(\approx k^3/2\) per product, light \(O(k^2)\) corrections) versus deep
single-product TA (\(\approx k^3/3\) per product, but \(\sim 4\times\) the
\(O(k^2)\) corrections).

At a small leaf the corrections dominate and pairs win decisively. For
the seven \(\nmpshape{7}{7}{7}\) leaves of a \(\nmpshape{14}{14}{14}\)
(or padded \(\nmpshape{17}{17}{17}\)) recombination:

\begin{center}
\begin{tabular}{lrr}
\toprule
option & per product & total (7 leaves) \\
\midrule
solo (catalog SOTA) & \(249\) & \(1743\) \\
pair-fused & \(245\) & \(\mathbf{1719}\) \\
full single-product TA & \(390\) & \(2730\) \\
\bottomrule
\end{tabular}
\end{center}

\noindent Pair fusion (\(1719\)) beats both solo (\(1743\)) and full TA
(\(2730\)). The deep TA's \(\tfrac{1}{3}\) only overcomes its heavier
corrections at \emph{even} \(k \ge 18\); odd \(k\) carries the heavier
\(15k^2\) term and never wins in the practical range. This is why the
catalog pair-fuses small leaves and reserves full TA for large even ones.
The chooser \code{PairFusedRecombination.chooseBest} weighs all three
options (\code{SOLO} / \code{PAIR\_FUSED} / \code{FULL\_TA}) per leaf and
returns the cheapest.

\subsection{Disjoint-sum decomposition ($\tau$-theorem)}
\label{ssec:disjoint-sum}

A target tensor can decompose as a sum of smaller matmul sub-tensors
\emph{at the tensor level} -- not at a single-axis level -- when the
target admits a $\tau$-theorem-style decomposition. The construction
is constructive: each leg of the disjoint sum computes a smaller
matmul, and the legs combine via the additive structure of the
tensor. Schwartz--Zwecher 2025 \cite{schwartzzwecher2025} give explicit
constructive disjoint-sum decompositions for \emph{large} formats
(e.g.\ \(n = 44, 56, 60\)). We caution that several FMM
catalog~\cite{fmmlille} entries that \emph{look} like $\tau$/direct-sum
identities at smaller shapes (e.g.\ \(\nmpshape{17}{17}{17}\) at rank
\(2934\)) turned out, on inspection, to be plain recombinations of a
small base rather than genuine disjoint sums. (Our own
\(\nmpshape{17}{17}{17}=2930\) is likewise a Winograd recombination,
Section \ref{ssec:axis-flip}, not a disjoint sum.)

We stress that \emph{we do not currently leverage the $\tau$-theorem
anywhere}. The disjoint-sum schemes we hold (Schwartz--Zwecher and
those FMM/Perminov build this way) are \emph{imported as explicit
factor matrices}, not reconstructed from a $\tau$ decomposition: no
catalog entry carries a $\tau$ lineage, and every gap we initially
attributed to a $\tau$ mechanism turned out, on inspection, to close
(or not) by other means -- Kronecker, recombination, or a fortunate
cousin scheme. A \code{DisjointSum} lineage node (with optional
\code{taLegs} groups) is \emph{reserved} in the schema for this
provenance but is exercised by no scheme today, and we have no search
operator that constructs disjoint sums. Section
\ref{sec:openquestions} flags this as the largest open extension to
the search.

\subsection{Summary table of operators}

\begin{table}[h]
\centering
\caption{Derivation strategies used by the catalog. ``Sharing'' is
the operator's ability to produce a scheme with fewer than
\(\sum_k r_{\text{sub-}k}\) bilinear products -- see Section
\ref{sec:nonoverlap}.}
\label{tab:strategies}
\begin{tabular}{lll}
\toprule
Operator & Source & Sharing? \\
\midrule
Axis-flip orbit & folklore + this paper \S\ref{ssec:axis-flip} & no \\
Kronecker product & Strassen 1969 & no \\
Axis concatenation & folklore & no \\
Recombination with allocation & Strassen-style block split & no \\
Output peel & DIS09 / Islam 2009 \(\gamma_5\) & no \\
Pair fusion (Pan TA) & Pan 1980 & limited (paired only) \\
Disjoint-sum / $\tau$\textsuperscript{$\dagger$} & Pan 1980 / Sch\"onhage 1981 / SZ 2025 & n/a (imported) \\
\bottomrule
\end{tabular}

\smallskip
\noindent\textsuperscript{$\dagger$}\,Not a search operator: disjoint-sum
schemes are imported as explicit factor matrices; we do not construct
them (Section~\ref{sec:openquestions}).
\end{table}

The ``no-sharing'' rows are the ones that fall under the non-overlap
property of Section \ref{sec:nonoverlap}. The ``yes-sharing'' rows
are deliberately written as separate constructors so that the
non-overlap invariant holds for every operator the generic
recombination path can produce.
